\ExplSyntaxOn

\int_new:N \l__binp_int
\int_set:Nn \l__binp_int {4}
\tl_new:N \l__bina_tl
\tl_new:N \l__binb_tl
\tl_new:N \l__binhw_tl
\int_new:N \l__bina_int
\int_new:N \l__binb_int
\int_new:N \l__binc_int
\int_new:N \l__bind_int
\fp_new:N \l__bina_fp

\DeclareDocumentCommand \BinaryPrecision {m}
{
  \int_set:Nn \l__binp_int {#1}
}

\cs_new_nopar:Npn \int_to_bin:Nn #1#2
{
  \tl_clear:N #1
  \int_set:Nn \l__bina_int {#2}
  \prg_replicate:nn {\l__binp_int}
  {
    \int_set:Nn \l__binb_int {\int_mod:nn {\l__bina_int} {2}}
    \tl_put_left:Nx #1 {\int_use:N \l__binb_int}
    \int_set:Nn \l__bina_int {\int_div_truncate:nn {\l__bina_int} {2}}
  }
}

\cs_new_nopar:Npn \bin_hamming_weight:NN #1#2
{
  \tl_set_eq:NN \l__binhw_tl #2
  \tl_replace_all:Nnn \l__binhw_tl {0} {}
  \int_set:Nn #1 {\tl_count:N \l__binhw_tl}
}

\cs_new_nopar:Npn \bin_flip_one:NNn #1#2#3
{
  \tl_set_eq:NN #2#1
  \prg_replicate:nn {#3 - 1}
  {
    \tl_replace_once:Nnn #2 {1} {a}
  }
  \tl_replace_once:Nnn #2 {1} {0}
  \tl_replace_all:Nnn #2 {a} {1}
}

\DeclareDocumentCommand \HammingWeight {m m}
{
  \int_to_bin:Nn \l__binhw_tl {#2}
  \bin_hamming_weight:NN #1 \l__binhw_tl
}

\DeclareDocumentCommand \HyperCube {m}
{
  \int_set:Nn \l__binp_int {#1}
  \int_zero_new:c {l__hyper_0_int}
  \int_set:cn {l__hyper_0_int} {-1}
  \int_step_inline:nnnn {1} {1} {#1}
  {
    \int_zero_new:c {l__hyper_##1_int}
    \int_set:cn  {l__hyper_##1_int} {\int_use:c {l__hyper_\int_eval:n {##1 - 1}_int} * (#1 - ##1 + 1) / ##1}
  }
  \fp_set:Nn \l__bina_fp {2^{#1}-1}
  \int_set:Nn \l__binc_int {\fp_to_int:N \l__bina_fp}
  \int_step_inline:nnnn {0} {1} {\l__binc_int}
  {
    \int_to_bin:Nn \l__bina_tl {##1}
    \bin_hamming_weight:NN \l__binb_int \l__bina_tl
    \node[every~ hypercube~ label/.try] (hg- \l__bina_tl) at (\int_use:c {l__hyper_\int_use:N \l__binb_int _int},\int_use:N \l__binb_int) {\texttt{\l__bina_tl}};
    \int_incr:c {l__hyper_\int_use:N \l__binb_int _int}
    \int_incr:c {l__hyper_\int_use:N \l__binb_int _int}
  }
  \int_step_inline:nnnn {0} {1} {\l__binc_int}
    {
      \int_to_bin:Nn \l__bina_tl {##1}
      \bin_hamming_weight:NN \l__binb_int \l__bina_tl
      \int_step_inline:nnnn {1} {1} {\l__binp_int}
      {
        \bin_flip_one:NNn \l__bina_tl \l__binb_tl {####1}
        \tl_if_eq:NNTF \l__bina_tl \l__binb_tl
          { % do nothing if nodes are the same
          }
          {
            \draw[every~ hypercube~ edge/.try] (hg- \l__bina_tl) -- (hg- \l__binb_tl);
          }
      }
    }
  }

\ExplSyntaxOff

% Macro to cross out a single node
\newcommand{\CrossOutNode}[1]{%
  \draw[red, thick]
    ($(#1) + (-0.3, 0.3)$)
    --
    ($(#1) + (0.3, -0.3)$);
  \draw[red, thick]
    ($(#1) + (-0.3, -0.3)$)
    --
    ($(#1) + (0.3, 0.3)$);
}

% Macro to cross out a list of nodes
\newcommand{\CrossOutNodes}[1]{%
  \foreach \n in {#1} {%
    \CrossOutNode{\n}%
  }
}
% Macro to draw a circle around a single node
\newcommand{\CircleNode}[2][RoyalBlue]{%
  \draw[#1, thick]
    ($(#2)$)
    circle (0.5);
}

% Macro to circle a list of nodes
\newcommand{\CircleNodes}[2][RoyalBlue]{%
  \foreach \n in {#2} {%
    \CircleNode[#1]{\n}%
  }
}

% command to make sure colorboxes don't create whitespace between code lines
\fboxsep=.6pt

% default settings for listings
\lstset{style=python}

\title[Datenintegrität]{Datenintegrität}

\subtitle{RAID}

\author[Blum, Cyril] % (optional)
{Cyril Blum}

\institute[KLW] % (optional)
{
    Fachschaft Informatik\\
    Kantonsschule im Lee
}

\date[\today] % (optional)

\logo{\includegraphics[width=2cm]{Logos/FreeFlower.pdf}}

\titlegraphic{\vspace{-.7cm}\includegraphics[width=.85\textwidth]{"Grundlagen_Info/06_Datenbanken/Figures/title-emojis"}}

\frame{\titlepage}

\begin{frame}{Problem: Bitrot}
    \begin{center}
        \includegraphics[width=.88\textwidth]{Grundlagen_Info/05_Datenintegritaet/Figures/bitrot.png}

        % Quelle: wahrscheinlichste Herkunft des Screenshots (nicht abschliessend bestätigt)
        {\tiny Quelle: Tony \& Chelsea Northrup, \enquote{NAS Storage for Photos \& Video: PROTECT YOUR DATA!}, YouTube, 2022, \url{https://www.youtube.com/watch?v=xvbnhCDfREk}}
    \end{center}
\end{frame}

\begin{frame}{Mögliche Lösung: \glsentryshort{NAS} mit \glsentryshort{RAID}}
    \begin{center}
        \includegraphics[height=.45\textheight]{Grundlagen_Info/05_Datenintegritaet/Figures/NAS.png}
    \end{center}
    \begin{itemize}
        \item \gls{NAS}: Netzwerkgebundener Speicher
        \item \gls{RAID}: Redundante Anordnung unabhängiger Festplatten
        \item \textbf{Achtung:} \acrshort{RAID} ist kein Backup! Gelöschte, überschriebene, verschlüsselte (Ransomware) oder beschädigte Dateien landen sofort auf allen Platten $\rightarrow$ separates Backup (offline / extern)
    \end{itemize}
\end{frame}


\begin{frame}{\glsentryshort{RAID} 0}
    \begin{center}
        \includegraphics[width=.17\textwidth]{Grundlagen_Info/05_Datenintegritaet/Figures/RAID_0.pdf}

        {\tiny Grafik: Cburnett, \href{https://commons.wikimedia.org/wiki/File:RAID_0.svg}{Wikimedia Commons}, CC BY-SA 3.0}
    \end{center}
    \begin{itemize}
        \item \textbf{Prinzip:} Striping (Daten werden in Blöcke aufgeteilt und auf die Festplatten verteilt)
        \item \textbf{Vorteil:} Hohe Geschwindigkeit
        \item \textbf{Nachteil:} Hohe Wahrscheinlichkeit für Datenverlust (kein Schutz)
        \item \textbf{Anwendung:} Temporäre / grosse Daten, z.B. Spiele
        \item \textbf{Kapazität:} $n$ Festplatten $\rightarrow$ $n$-fache Kapazität
    \end{itemize}
\end{frame}

\begin{frame}{\glsentryshort{RAID} 1}
    \begin{center}
        \includegraphics[width=.17\textwidth]{Grundlagen_Info/05_Datenintegritaet/Figures/RAID_1.pdf}

        {\tiny Grafik: Cburnett, \href{https://commons.wikimedia.org/wiki/File:RAID_1.svg}{Wikimedia Commons}, CC BY-SA 3.0}
    \end{center}
    \begin{itemize}
        \item \textbf{Prinzip:} Mirroring (je zwei Festplatten bilden ein identisches Spiegelpaar)
        \item \textbf{Vorteil:} Hohe Ausfallsicherheit (pro Spiegelpaar kann eine Festplatte ausfallen)
        \item \textbf{Nachteil:} Nur halbe nutzbare Kapazität
        \item \textbf{Anwendung:} Wichtige Daten, z.B. Server
        \item \textbf{Kapazität:} $n$ Festplatten $\rightarrow$ $n/2$-fache Kapazität
    \end{itemize}
\end{frame}

\begin{frame}{\glsentryshort{RAID} 5}
    \begin{center}
        \includegraphics[width=.36\textwidth]{Grundlagen_Info/05_Datenintegritaet/Figures/RAID_5.pdf}

        {\tiny Grafik: Cburnett, \href{https://commons.wikimedia.org/wiki/File:RAID_5.svg}{Wikimedia Commons}, CC BY-SA 3.0}
    \end{center}
    \begin{itemize}
        \item \textbf{Prinzip:} Striping mit verteilter Parität (Paritätsinformationen werden auf alle Festplatten verteilt)
        \item \textbf{Vorteil:} Gute Geschwindigkeit, Schutz bei Ausfall einer Festplatte
        \item \textbf{Nachteil:} Wiederherstellung bei Ausfall dauert lange
        \item \textbf{Anwendung:} Dateiserver, \gls{NAS}
        \item \textbf{Kapazität:} $n$ Festplatten $\rightarrow$ $(n-1)$-fache Kapazität
    \end{itemize}
\end{frame}

\begin{frame}{\glsentryshort{RAID} 6}
    \begin{center}
        \includegraphics[width=.5\textwidth]{Grundlagen_Info/05_Datenintegritaet/Figures/RAID_6.pdf}

        {\tiny Grafik: Cburnett, \href{https://commons.wikimedia.org/wiki/File:RAID_6.svg}{Wikimedia Commons}, CC BY-SA 3.0}
    \end{center}
    \begin{itemize}
        \item \textbf{Prinzip:} Striping mit doppelter Parität (zwei Paritätsblöcke pro Stripe)
        \item \textbf{Vorteil:} Schutz bei Ausfall von zwei Festplatten
        \item \textbf{Nachteil:} Noch geringere nutzbare Kapazität, langsamere Schreibgeschwindigkeit
        \item \textbf{Anwendung:} Grosse Speicherlösungen mit hoher Sicherheit
        \item \textbf{Kapazität:} $n$ Festplatten $\rightarrow$ $(n-2)$-fache Kapazität
    \end{itemize}
\end{frame}

\begin{frame}{\glsentryshort{RAID} ist kein Backup}
    \begin{columns}[T,onlytextwidth]
        \column{.48\textwidth}
        \textbf{\acrshort{RAID} schützt vor \dots}
        \begin{itemize}
            \item Ausfall einer Festplatte (\acrshort{RAID} 6: zwei)
            \item Stillstand: das System läuft weiter
        \end{itemize}
        \column{.48\textwidth}
        \textbf{\acrshort{RAID} schützt \textit{nicht} vor \dots}
        \begin{itemize}[<+->]
            \item versehentlichem Löschen / Überschreiben
            \item Ransomware
            \item Diebstahl, Brand, Wasserschaden
            \item \textit{bit rot}, das unbemerkt bleibt
        \end{itemize}
    \end{columns}
    \vspace{4mm}
    \begin{center}
        \only<5>{Jede Änderung landet sofort auf \textit{allen} Festplatten $\rightarrow$ es braucht eine \textbf{separate} Kopie: ein \textbf{Backup}.}
    \end{center}
\end{frame}

\begin{frame}{Die 3-2-1-Regel}
    \begin{center}
        \resizebox{.95\textwidth}{!}{% 3-2-1-Backup-Regel: 3 Kopien, auf 2 verschiedenen Medien, davon 1 ausser Haus.
% Wird im Skript (Abschnitt Backup) und auf den RAID-Folien verwendet.
\begin{tikzpicture}[
        copy/.style={draw, rounded corners, thick, minimum width=3.2cm, minimum height=1.5cm, align=center, font=\small},
        zone/.style={draw, dashed, rounded corners, inner sep=8pt},
        lbl/.style={font=\footnotesize\bfseries, anchor=south west}
    ]
    % Zuhause: Original und erste Kopie auf einem anderen Medium
    \node[copy, fill=blue!10]   (orig)  at (0,0)   {\faLaptop\\[2pt] Original\\ \footnotesize Laptop (SSD)};
    \node[copy, fill=green!12]  (copy1) at (4,0)   {\faHdd\\[2pt] Kopie 1\\ \footnotesize externe Festplatte / \acrshort{NAS}};
    % Ausser Haus: zweite Kopie
    \node[copy, fill=orange!15] (copy2) at (9.8,0)   {\faCloud\\[2pt] Kopie 2\\ \footnotesize Cloud / Festplatte bei Verwandten};

    \node[zone, fit=(orig)(copy1)]  (home) {};
    \node[zone, fit=(copy2)]        (away) {};
    \node[lbl] at (home.north west) {Zuhause};
    \node[lbl] at (away.north west) {Ausser Haus};

    \draw[-latex, thick] (orig) -- (copy1);
    \draw[-latex, thick] (copy1) -- (copy2);

    % Die drei Regeln unter der Grafik
    \node[font=\small, align=center] at (4.9,-1.9)
    {\textbf{3} Kopien \quad\textperiodcentered\quad auf \textbf{2} verschiedenen Medien \quad\textperiodcentered\quad davon \textbf{1} ausser Haus};
\end{tikzpicture}
}
    \end{center}
    \begin{itemize}
        \item Backups automatisch und regelmässig erstellen, ältere Versionen aufbewahren
        \item Wiederherstellung testen -- erst dann weiss man, dass das Backup funktioniert
    \end{itemize}
\end{frame}

\begin{frame}{Auftrag}{Skript \enquote{Datenintegrität}}
    \begin{itemize}
        \item \aufgabebox{1.27--1.28 (\acrshort{XOR})}
        \item \aufgabebox{1.29--1.34 (\acrshort{RAID})}
        \item \aufgabebox{1.35--1.36 (Backup)}
    \end{itemize}
\end{frame}